\section{Discussion}

\subsection{Factors Contributing to Performance Improvement by Depth Guidance}
In this study, methods incorporating depth information—namely Depth-guided Attention and Depth-guided Augmentation—consistently achieved higher classification accuracy than ResNet18 \cite{ref12} and ViT \cite{ref1}. These results indicate that depth guidance suppresses the influence of background regions and directs model attention toward foreground structures such as trunks, branches, and leaves. In particular, Depth-guided Augmentation enables simulation of diverse environmental conditions through background replacement, thereby promoting robust feature learning that is less dependent on background context. This effect was also supported by qualitative visualization results, in which depth-guided models exhibited stronger focus on the foreground.

\subsection{Influence of DINO Pretraining}
Although Depth-guided Attention showed smaller improvement margins compared with augmentation, it maintained stable foreground focus during inference, likely because the depth weighting was represented as a continuous rather than discrete signal throughout training. This suggests that integrating depth directly into the model architecture—rather than relying solely on data augmentation—offers a more generalizable and deployable approach in real-world applications.

Self-supervised pretraining with DINO \cite{ref13} enables the model to learn global structural representations from large-scale natural images without label supervision. This property allows the model to capture abstract shape and contextual relations within a high-dimensional feature space. However, DINO tends to prioritize global scene consistency over local object details; hence, when applied to tree images where background and object regions are highly entangled, attention often spreads across background areas. Consequently, configurations without depth guidance showed decreased accuracy.

In contrast, when depth guidance was combined with DINO, performance improved in most configurations. This complementary effect can be interpreted as follows: DINO’s global representation captures holistic scene structure, while depth information acts as a local suppression signal that reweights features toward the foreground. In other words, depth guidance constrains DINO’s abstract, scene-level representations into more object-centered ones, allowing hierarchical features such as trunks, branches, and leaves to be extracted more distinctly. Thus, the integration of DINO and depth guidance achieves robust discrimination even under background noise.

\subsection{Stability of Late Fusion and Piecewise Depth}
Introducing Late Fusion consistently degraded accuracy across all settings. The logarithmic combination of depth weights alters the relative strength of attention scores, causing the attention distribution to either overexpand or become locally biased. Such score fluctuations hinder consistent focus on major foreground regions. Attention map visualizations also revealed attention drifting toward the ground and background, suggesting that simple additive fusion disrupts the scale consistency between depth and feature representations.

On the other hand, Piecewise Depth produced limited improvements for the UrbanStreetTree dataset but decreased accuracy for FruitsPark. This discrepancy likely stems from differences in object distance and background composition between datasets, meaning a single discretization threshold for depth cannot be universally applied. Therefore, to make the piecewise approach effective, dynamic segmentation or learnable optimization based on depth distribution is necessary.

\subsection{Dataset Characteristics and Generalization}
While all depth-guided methods improved accuracy on the UrbanStreetTree dataset, several configurations resulted in degraded performance on FruitsPark. This can be attributed to the diverse illumination conditions, seasonal variations, and the presence of nearby ground and neighboring trees in FruitsPark, which introduce strong visual noise. Such environmental complexity can cause the depth estimation model to misclassify foreground and background regions, thereby reducing the reliability of depth-based attention guidance. Consequently, the effectiveness of depth guidance strongly depends on the accuracy of depth estimation and the stability of environmental imaging conditions.

